\documentclass[11pt]{article}
\usepackage{docstyle}
\renewcommand\sheetname{Handout 2 \textperiodcentered{} Exponentials and logarithms}

\begin{document}
\sheethead{Lesson 2 \textperiodcentered{} Exponentials and Logarithms}
  {NCEA Level 2 Mathematics \textperiodcentered{} Algebra}{Handout}

\begin{keybox}[title={The one idea: a log is a power}]
$\log_{b} y$ asks one question: \textbf{what power of $b$ gives $y$?}
\[ 2^{3} = 8 \quad\Longleftrightarrow\quad \log_{2} 8 = 3 \qquad\qquad y = b^{x} \quad\Longleftrightarrow\quad \log_{b} y = x \]
Every question in this lesson is a switch between these two forms. Say it as you write it: \textbf{base to the bottom, $y$ in the brackets}.
\end{keybox}

% =================================================================== 1
\section{What a log means}\label{h2:meaning}
\begin{minipage}[t]{0.48\textwidth}\vspace{0pt}
\renewcommand\arraystretch{1.5}
\begin{tabular}{@{}l l l@{}}
\toprule
\textbf{Power form} & \textbf{Log form} & \textbf{Value}\\ \midrule
$3^{4} = 81$ & $\log_{3} 81 = 4$ & 4 \\
$10^{-2} = 0.01$ & $\log 0.01 = -2$ & $-2$ \\
$9^{\frac{1}{2}} = 3$ & $\log_{9} 3 = \frac{1}{2}$ & $\frac{1}{2}$ \\
$b^{0} = 1$ & $\log_{b} 1 = 0$ & 0 \\
$b^{1} = b$ & $\log_{b} b = 1$ & 1 \\
\bottomrule
\end{tabular}
\end{minipage}\hfill
\begin{minipage}[t]{0.48\textwidth}\vspace{0pt}
\begin{itemize}
\item $\log$ with no base written means $\log_{10}$: the calculator's \textbf{log} key.
\item $\ln$ means $\log_{e}$, where $e \approx 2.718$: the calculator's \textbf{ln} key. $e^{x}$ and $\ln x$ undo each other.
\item The \Tbase{} is positive and not 1. You can only take the log of a \textbf{positive} number.
\end{itemize}
\end{minipage}

\subsection*{Write it in this order (an equation with one log)}
\begin{enumerate}
\item Get the log on its own: $\log_{b}(\text{something}) = k$.
\item Switch to power form: $\text{something} = b^{k}$.
\item Solve. If the unknown is the base, the base must be positive: $\log_{x} 49 = 2$ gives $x^{2} = 49$, so $x = 7$, not $-7$.
\end{enumerate}

\textbf{Example 1.} Solve $\log_{5} x = 3$.\quad $x = 5^{3} = \mathbf{125}$ \mtag{u}

% =================================================================== 2
\section{The log laws}\label{h2:laws}
These are on the formula sheet. Each one is an index law in disguise.
\begin{center}
\renewcommand\arraystretch{1.6}
\begin{tabular}{@{}l l l@{}}
\toprule
\textbf{Log law} & \textbf{Index law behind it} & \textbf{Example} \\ \midrule
$\log_{b}(xy) = \log_{b}x + \log_{b}y$ & multiply: add the powers & $\log_{2}3 + \log_{2}5 = \log_{2}15$ \\
$\log_{b}\Bigl(\dfrac{x}{y}\Bigr) = \log_{b}x - \log_{b}y$ & divide: subtract the powers & $\log_{2}12 - \log_{2}4 = \log_{2}3$ \\
$\log_{b}(x^{n}) = n\log_{b}x$ & power of a power: multiply & $\log x^{3} = 3\log x$ \\
\bottomrule
\end{tabular}
\end{center}
\textbf{There is no law for} $\log(x+y)$, and $\dfrac{\log x}{\log y}$ is \textbf{not} $\log\dfrac{x}{y}$.

\subsection*{Write it in this order (a single log)}
\begin{enumerate}
\item Move every number in front of a log up as a power: $2\log x = \log x^{2}$.
\item Plus signs multiply inside one log; minus signs divide.
\item Simplify the fraction inside.
\end{enumerate}

\begin{multicols}{2}
\textbf{Example 2.} Write $2\log x + \log 5 - \log 10$ as a single log.
\[ = \log x^{2} + \log 5 - \log 10 = \log\frac{5x^{2}}{10} = \mathbf{\log\frac{x^{2}}{2}} \]
\mtag{u}
\columnbreak

\textbf{Example 3.} Solve $\log_{3}(x+4) + \log_{3}(x-4) = 2$.
\begin{align*}
\log_{3}\bigl((x+4)(x-4)\bigr) &= 2\\
x^{2} - 16 &= 3^{2} = 9\\
x &= 5 \quad (x = -5 \text{ rejected})
\end{align*}
\mtag{r}
$x = -5$ makes $x - 4$ negative, and the log of a negative number does not exist.
\end{multicols}

% =================================================================== 3
\section{Solving exponential equations}\label{h2:solve}
\begin{keybox}[title={Write it in this order (unknown in the power)}]
\begin{enumerate}
\item Get the power on its own: $b^{(\ldots)} = k$.
\item Same base on both sides? Match the powers and solve exactly. Otherwise:
\item Take logs of both sides and bring the power down \textbf{in a bracket}: $(2x-1)\log 5 = \log 200$.
\item Divide by the log, then finish the rearranging. Round only at the end.
\item For $e$, use $\ln$: $e^{0.3t} = 5$ gives $0.3t = \ln 5$.
\end{enumerate}
\end{keybox}

\begin{multicols}{2}
\textbf{Example 4.} Solve $5^{2x-1} = 200$.
\begin{align*}
(2x-1)\log 5 &= \log 200\\
2x - 1 &= \frac{\log 200}{\log 5} = 3.2920\ldots\\
x &= \mathbf{2.146}
\end{align*}
\mtag{u}

\textbf{Example 5.} Solve $4^{x+1} = 8^{x}$ exactly.
\[ 2^{2(x+1)} = 2^{3x} \;\Rightarrow\; 2x + 2 = 3x \;\Rightarrow\; x = \mathbf{2} \]
\mtag{u}
\columnbreak

\textbf{Example 6.} Solve $3^{2x} - 4\times 3^{x} + 3 = 0$.\par
Let $u = 3^{x}$, so $3^{2x} = u^{2}$:
\begin{align*}
u^{2} - 4u + 3 &= 0\\
(u-1)(u-3) &= 0\\
3^{x} = 1 \text{ or } 3^{x} &= 3\\
x = \mathbf{0} \text{ or } x &= \mathbf{1}
\end{align*}
\mtag{t}
\end{multicols}

% =================================================================== 4
\section{Growth and decay}\label{h2:growth}
\[ P = A\,r^{t} \qquad A = \text{starting amount},\quad r = \text{\Tgrowth{} per time step},\quad t = \text{number of steps} \]
\begin{itemize}
\item Growth of $R\%$: $r = 1 + \dfrac{R}{100}$ (6\% growth: $r = 1.06$). Decay of $R\%$: $r = 1 - \dfrac{R}{100}$ (15\% loss: $r = 0.85$).
\item \textbf{When} questions: put the numbers in, divide by $A$, take logs. \textbf{What rate} questions: divide, then take the root ($t$-th root).
\item Finish with \textbf{a sentence in context}, with units, and say what the decimal means (``during the 7th year'').
\end{itemize}

\begin{multicols}{2}
\textbf{Example 7.} A car bought for \$18\,000 loses 15\% of its value each year. When is it first worth less than \$6000?
\begin{align*}
18\,000(0.85)^{t} &= 6000\\
0.85^{t} &= \tfrac{1}{3}\\
t &= \frac{\log\frac{1}{3}}{\log 0.85} = 6.76
\end{align*}
\mtag{r}
It is first worth less than \$6000 \textbf{during the 7th year}, after about 6.8 years.
\columnbreak

\textbf{Example 8.} \$5000 grows to \$7500 in 8 years at a fixed annual rate $R\%$. Find $R$.
\begin{align*}
5000\Bigl(1 + \frac{R}{100}\Bigr)^{8} &= 7500\\
1 + \frac{R}{100} &= 1.5^{\frac{1}{8}} = 1.0520\\
R &= \mathbf{5.2\%}
\end{align*}
\mtag{r}
\end{multicols}

% =================================================================== 5
\section{Versions the marker won't accept}\label{h2:wontaccept}
\begin{tabularx}{\textwidth}{@{}>{\raggedright}p{0.3\textwidth} >{\raggedright}p{0.24\textwidth} X@{}}
\toprule
\textbf{Question} & \textbf{Answer given} & \textbf{Why it falls short} \\ \midrule
Show that $D = \dfrac{\log 2}{\log\bigl(1 + \frac{R}{100}\bigr)}$ & the final line only & A ``show that'' needs the line where logs are taken of both sides: markers look for that line. \tabularnewline
Solve $\log_{x} 49 = 2$ & $x = \pm 7$ & A base cannot be negative. The answer is $x = 7$ with $-7$ rejected for that reason. \tabularnewline
How long until \ldots? & $t = 6.7599$ & No units, no meaning. Give a sentence: ``during the 7th year''. \tabularnewline
Write as a single logarithm & $\log 3a + \log\dfrac{a^{2}}{36}$ & Two logs are not a single log, and the fraction inside is not simplified. \tabularnewline
\bottomrule
\end{tabularx}

% =================================================================== 6
\section{Ways to lose marks}\label{h2:lose}
From the wording of the marking schedules and from teaching observation; numbered for marking.
\begin{description}[leftmargin=10mm, labelwidth=9mm, labelsep=1mm, font=\normalfont]
\item[\textbf{\Lref{logform}}] \textbf{\Ltext{logform}.} $\log_{2} 32 = 5$ means $2^{5} = 32$. Base to the bottom, $y$ in the brackets.
\item[\textbf{\Lref{loglaw}}] \textbf{\Ltext{loglaw}.} $\log(x+3) \neq \log x + \log 3$, and $\dfrac{\log 40}{\log 5} \neq \log 8$.
\item[\textbf{\Lref{bracketpower}}] \textbf{\Ltext{bracketpower}.} $\log 5^{x-2} = (x-2)\log 5$, not $x - 2\log 5$.
\item[\textbf{\Lref{argument}}] \textbf{\Ltext{argument}.} Check every answer in the original logs: each bracket must be positive, and a base must be positive.
\item[\textbf{\Lref{factor}}] \textbf{\Ltext{factor}.} A 20\% loss each year multiplies by $0.8$, not $0.2$.
\item[\textbf{\Lref{context}}] \textbf{\Ltext{context}.} Units, a sentence, and rounding only at the end.
\end{description}

% =================================================================== 7
\section{Quick review}\label{h2:review}
\renewcommand\arraystretch{1.45}
\begin{minipage}[t]{0.47\textwidth}\vspace{0pt}
\textbf{Powers to know by sight}\par\smallskip
\begin{tabular}{@{}l l l l@{}}
\toprule
$2^{5} = 32$ & $2^{6} = 64$ & $2^{7} = 128$ & $2^{10} = 1024$\\
$3^{4} = 81$ & $3^{5} = 243$ & $5^{3} = 125$ & $5^{4} = 625$\\
$4^{3} = 64$ & $6^{3} = 216$ & $7^{3} = 343$ & $10^{-2} = 0.01$\\
\bottomrule
\end{tabular}
\end{minipage}\hfill
\begin{minipage}[t]{0.47\textwidth}\vspace{0pt}
\textbf{Index laws behind the log laws}\par\smallskip
\begin{tabular}{@{}>{\raggedright}p{0.42\linewidth} >{\raggedright\arraybackslash}p{0.5\linewidth}@{}}
\toprule
$x^{a}\times x^{b} = x^{a+b}$ & $\log(xy) = \log x + \log y$ \\
$x^{a} \div x^{b} = x^{a-b}$ & $\log\frac{x}{y} = \log x - \log y$ \\
$(x^{a})^{n} = x^{an}$ & $\log x^{n} = n\log x$ \\
\bottomrule
\end{tabular}
\end{minipage}
\end{document}
